\documentclass[12pt]{article}
\usepackage[utf8]{inputenc}
\usepackage{setspace}
\usepackage{graphicx}
\usepackage{amsmath}
\usepackage{booktabs}
\usepackage{float}
\usepackage[margin=1in]{geometry}
\usepackage{indentfirst} 
\usepackage{natbib}

\setstretch{1.5} 
\setlength{\parindent}{0.5in} 

\title{Liverpool F.C. Performance Analysis\\
A Comparison of the 2023-24 and 2024-25 Seasons}
\author{Jack Vincent}
\date{May 7, 2025}

\begin{document}

\maketitle

\section*{Introduction}

What factors drive a football club's championship-winning season? The examination of key performance indicators influencing match outcomes has profound implications for club strategy and tactical approaches. I've taken on this question by examining Liverpool Football Club's contrasting performances across two Premier League seasons: their third-place finish in 2023-24 under Jürgen Klopp and their championship victory in 2024-25 under new manager Arne Slot.

This transition from Klopp to Slot marked a critical turning point for Liverpool. Many wondered if the club could maintain competitive performance during such a significant leadership change. The remarkable improvement under Slot presents an ideal case study to identify which performance metrics truly drive championship-level results in elite football.

My investigation centers on three distinct performance indicators: Expected Goals Difference (xG Difference), Shot Conversion Rate, and Home Advantage. Expected Goals has revolutionized football analytics by providing a sophisticated measurement of chance quality rather than simply counting shots \citep{cavus2022}. By calculating the gap between expected goals created versus conceded, xG Difference effectively combines offensive creation with defensive solidity into a single metric. Shot Conversion Rate offers insight into finishing efficiency independent of chance quality, potentially revealing tactical improvements in attacking effectiveness. Home Advantage serves as a control variable, helping isolate performance changes across different managerial approaches.

Through regression analysis of match-level data, I've identified which factors most significantly influenced Liverpool's results across both seasons. Rather than relying on aggregate season statistics, this approach captures the nuanced performance patterns that contributed to their championship.

This research advances our understanding of football performance analysis by offering a data-driven assessment of factors facilitating successful managerial transitions. The findings provide practical guidance for clubs facing leadership changes and demonstrate how advanced metrics can reveal insights beyond traditional statistics, creating a more comprehensive framework for evaluating on-field effectiveness.

\section*{Literature Review}

\subsection*{Expected Goals and Performance Measurement}

Expected goals represents a breakthrough in football analytics, offering a more nuanced evaluation of team and player performance than traditional metrics. \citet{cavus2022}'s analysis of over 315,000 shots across seven seasons in Europe's top leagues reveals that xG provides deeper insights than conventional statistics like shots on target or possession percentage. They note that xG ``provides a more representative measure of the team and player performance which also suit the low-scoring nature of football'' (p. 2). This statistical approach assigns probability values to shots based on factors like location, angle, and defensive pressure, creating a more sophisticated understanding of offensive effectiveness.

Building on this foundation, \citet{deb2019} developed a spatial modeling approach for shot conversion incorporating positional and situational variables. Their research demonstrates how conversion probability varies based on shot location, body part used, and defensive pressure. Their model ``identifies the factors significantly affecting the shot conversion rate and provides an estimate of the probability of scoring a goal if the location and other details are given'' (p. 4). This spatial perspective enriches our understanding of shooting efficiency, providing valuable context for analyzing Liverpool's performance across different managerial regimes.

\subsection*{Predictive Power of Advanced Metrics}

Beyond measuring individual shot quality, xG difference captures the balance of opportunity creation between teams. \citet{roccetti2024} evaluated the relationship between xG metrics and actual results using matches from the 2020 UEFA European Football Championship. Their analysis confirms that xG demonstrates ``in the long run, a good predictive power'' for match outcomes, validating my use of xG difference as a key variable.

The xG difference concept offers a more holistic performance measure than examining offensive or defensive metrics in isolation. By calculating the gap between expected goals created and conceded, we can evaluate both attacking productivity and defensive resilience simultaneously. Roccetti's team employed Kolmogorov-Smirnov hypothesis testing to demonstrate no statistically significant difference between tournament rankings based on actual goals versus xG metrics, further validating xG-based evaluation.

Efficiency measurement adds another valuable perspective to performance analysis. \citet{kern2012} proposed a two-stage data envelopment approach for the English Premier League, separating on-field from off-field efficiency. They argue this approach ``can identify sources of inefficiencies more accurately than a one-stage data envelopment analysis and provides football officials with valuable information about their club'' (p. 3). This distinction proves particularly relevant when examining Liverpool's transition between managers, as it helps distinguish between tactical effectiveness and resource utilization.

\subsection*{Team Performance and Contextual Factors}

Research specific to Liverpool FC provides important background context. \citet{isik2020} analyzed Liverpool's performance against top-6 opponents and in derby matches, examining whether player performances differ in high-stakes fixtures. By dividing players into positional groups and calculating z-scores across performance metrics, Işık created a template for comparative analysis across different match contexts.

The relationship between expected goals and actual scoring represents a critical area of investigation. While xG predicts scoring probability from various situations, conversion rates vary based on factors beyond shot location. \citet{anzer2023} addressed this gap by improving xG models to account for player quality and psychological factors. They note that traditional xG models ``have not until this point considered using important features, e.g., player/team ability and psychological effects'' (p. 1). This insight highlights the need to examine both chance quality through xG metrics and the efficiency with which teams convert opportunities into goals.

Home advantage consistently emerges as a significant factor in football outcomes. \citet{kern2012} acknowledge the importance of venue in their efficiency analysis, recognizing that performance expectations must account for the established advantage of playing at home. The home advantage factor becomes particularly interesting when comparing managerial approaches, as changes in home versus away form may reflect differences in tactical preparation or team mentality.

The most insightful approaches to football performance analysis integrate multiple perspectives. \citet{deb2019} exemplify this by combining spatial modeling with traditional statistical analysis to gain deeper insights into goal-scoring probability. This integrated view informs my analytical framework -- by examining both expected goals difference and shot conversion rate alongside categorical factors like venue and season, I acknowledge football performance's multi-dimensional nature.

Methodological considerations in football analytics vary considerably. \citet{isik2020} highlights the importance of normalizing performance metrics to account for positional differences and contextual factors. \citet{cavus2022} address the challenges of working with unbalanced datasets in football analytics, noting that most shots don't result in goals. These perspectives have shaped my analytical choices, particularly the decision to center variables and focus on relative performance differences rather than absolute values.

\section*{Data}

The primary data source for this project was gathered from FBref.com

For this analysis, I compiled match-level data for Liverpool FC across the 2023-24 and 2024-25 Premier League seasons. I gathered this information through web scraping and the worldfootballR package in R, which provides access to statistics from fbref.com. Using functions like fb\_team\_match\_results() and fb\_team\_match\_log\_stats(), I extracted consistent team performance metrics across both seasons.

My dataset comprises Liverpool's Premier League matches through May 2025, with the 2023-24 season limited to the first 34 matches to ensure comparability with the available 2024-25 data. For each match, I extracted:

\begin{enumerate}
\item Match Information: date, opponent, venue, and result
\item Shot Statistics: total shots, goals scored, and expected goals for both Liverpool and opponents
\item Performance Metrics: derived variables including xG difference and shot conversion rate
\item Categorical Variables: season indicator (0 for 2023-24, 1 for 2024-25) and home game indicator (1 for home, 0 for away)
\end{enumerate}

The expected goals data comes from the xG model implemented by Statsbomb and reported through fbref.com. This model estimates goal probability for each shot based on factors like shot location, body part used, pass type, and defensive pressure, providing a sophisticated measure of chance quality beyond raw goal counts.

After gathering the data, I calculated the dependent variable, PointsPerGame, coded as 3 for wins, 1 for draws, and 0 for losses. The final dataset contains complete information for all analyzed matches with no missing values for key variables, allowing for robust statistical analysis of performance metrics across the two managerial approaches.

\section*{Methods}

I've employed multiple linear regression to identify the key performance factors behind Liverpool's improved results under Arne Slot. This analytical approach allows me to quantify relationships between various performance indicators and match outcomes while controlling for potential confounding variables. Similar techniques have proven effective in previous studies examining connections between advanced metrics and results \citep{anzer2023}.

\subsection*{Regression Model Design}

My core analysis centers on a multiple linear regression model with points earned per match as the dependent variable. I've designed the model to include both main effects and interaction terms, capturing direct influences of performance metrics and their potentially varying impacts across seasons:
\begin{align}
\text{PointsPerGame} = \beta_0 &+ \beta_1(\text{xG\_Difference\_c}) + \beta_2(\text{ShotConversionRate\_c}) \notag \\
&+ \beta_3(\text{HomeGame}) + \beta_4(\text{Season}) \notag \\
&+ \beta_5(\text{xG\_Difference\_c} \times \text{Season}) \notag \\
&+ \beta_6(\text{ShotConversionRate\_c} \times \text{Season}) + \varepsilon
\end{align}

In this model, PointsPerGame represents match outcome (3 for wins, 1 for draw, 0 for losses); xG\_difference\_c is the centered Expected Goals Difference; ShotConversionRate\_c is the centered Shot Conversion Rate; HomeGame indicates venue (1 for home, 0 for away); and Season distinguishes between campaigns (0 for 2023-24, 1 for 2024-25). The interaction terms capture changes in how performance metrics influenced outcomes across seasons, with the error term $\varepsilon$ accounting for unexplained variance.

The inclusion of interaction terms builds on \citet{isik2020}'s finding that performance metrics may impact outcomes differently across varying match contexts. These terms allow me to formally test whether relationships between performance indicators and results differed significantly under Klopp versus Slot.

\subsection*{Variable Treatment and Diagnostic Testing}

To address potential multicollinearity concerns, particularly with interaction terms included, I centered the continuous predictors by subtracting their means. This technique reduces correlation between main effects and interaction terms while maintaining the original data's variability \citep{roccetti2024}. The transformed variables (xG\_Difference\_c and ShotConversionRate\_c) have means of zero but preserve the pattern of relationships in the original dataset.

After model estimation, I conducted thorough multicollinearity diagnostics using Variance Inflation Factors. This approach quantifies the extent to which coefficient variances are inflated due to correlations between predictors. All VIF values remained below 2.5, well below the commonly accepted threshold of 10, indicating minimal multicollinearity impact on coefficient precision \citep{obrien2007}.

\subsection*{Analytical Process}

I performed the regression analysis in R using the base lm() function for model fitting and the car package for variance inflation factor calculation. My analytical process began with exploratory data visualization to identify potential relationships and outliers, followed by model specification including both main effects and interaction terms. After estimation using ordinary least squares regression, I conducted diagnostic testing for multicollinearity using VIF values. Finally, I interpreted coefficient estimates, focusing on both statistical significance (p $<$ 0.05) and practical importance.

This approach balances statistical rigor with practical interpretability, providing robust examination of which performance metrics most influenced Liverpool's outcomes and how these relationships evolved across seasons under different managers.

\subsection*{Graphical Visualizations}

To translate my statistical findings into visually compelling evidence, I've created seven figures that capture different dimensions of Liverpool's performance evolution. Figure 1 presents the practical impact of key performance metrics on points earned, scaled to show real-world effects. This bar chart reveals shot conversion rate's dominant influence; a 10\% improvement generates approximately 1.12 additional points per match, substantially outweighing the impact of both home advantage and xG difference.

In Figure 2, I've compared average points earned by season and location, uncovering Slot's remarkable achievement in maintaining consistent performance across both home and away fixtures. This pattern contrasts sharply with the previous season's more variable results under Klopp and helps explain Liverpool's championship success despite similar statistical fundamentals.

The forest plot in Figure 3 displays all regression coefficients with their 95\% confidence intervals, allowing me to visually assess both statistical significance and effect size simultaneously. This visualization immediately highlights which factors significantly influenced match outcomes and which showed meaningful evolution between managerial approaches.

I found the interaction between shot conversion rate and season particularly intriguing, so Figure 4 specifically examines this relationship. The diverging lines visually confirm the significant negative interaction term from my regression model, showing how finishing efficiency's importance diminished under Slot's system while remaining critically important under Klopp's approach.

To further investigate venue effects, Figure 5 breaks down home versus away performance patterns across seasons. This more detailed look strengthens my finding that Slot's system produced more adaptable performance across different match contexts compared to the previous campaign.

For model validation, Figure 6 plots observed versus predicted points, helping me identify matches where my model performed particularly well or struggled. This diagnostic visualization validates the overall modeling approach while highlighting specific fixtures that might warrant closer tactical analysis.

Finally, Figure 7's residual analysis examines error patterns across seasons, confirming no systematic bias in my model and supporting the robustness of my statistical findings. These residual patterns help identify potential areas for model refinement in future research.

Together, these visualizations transform complex statistical relationships into clear visual evidence. They reveal not just which metrics drove Liverpool's success, but how the relationship between these metrics and match outcomes evolved during the transition from Klopp to Slot, demonstrating the value of this analytical approach for understanding managerial transitions in elite football.

\section*{Findings}

\subsection*{Regression Analysis Results}

The main results are reported in Table 2.

My linear regression analysis revealed several significant patterns in Liverpool's performance across the two seasons. The model achieved moderate explanatory power with an adjusted R-squared of 0.3718, indicating that approximately 37\% of the variation in points earned could be explained by the included predictors. While substantial variance remains unexplained, the highly significant F-statistic (p $<$ 0.001) confirms the model provides meaningful explanatory value.

Multicollinearity diagnostics confirmed the absence of problematic correlations among predictors, with all VIF values well below the accepted threshold of 10. The highest value of 2.4344 for ShotConversionRate\_c indicates minimal inflation of standard errors due to multicollinearity.

\subsection*{Analysis of Metrics}

The most remarkable finding was the substantial impact of shot conversion rate on match outcomes. With a coefficient of 11.1988 (p $<$ 0.001), finishing efficiency emerged as the strongest predictor of points earned. This coefficient suggests that a 100\% increase in conversion efficiency would yield an additional 11.2 points on average, holding other factors constant. In more practical terms, a 10\% improvement in Liverpool's ability to convert shots into goals would generate approximately 1.12 additional points per match, providing a meaningful advantage in tight championship races.

However, I discovered an unexpected pattern in the significant negative interaction between shot conversion rate and season (-8.0318, p $<$ 0.01). This interaction reveals that the relationship between conversion efficiency and points weakened substantially in the 2024-25 season compared to the previous campaign. This finding suggests that while finishing efficiency remained important under both managers, other factors played a more prominent role in determining outcomes during the championship season under Slot.

Expected goals difference also significantly predicted match outcomes, with a coefficient of 0.2148 (p $<$ 0.05). This indicates that for each additional expected goal Liverpool created above their opponent, they gained approximately 0.21 more points on average. While statistically significant, this effect is notably smaller than that of shot conversion rate, suggesting that finishing efficiency may have been more crucial than the raw quantity or quality of chances created.

The interaction between xG difference and season was positive (0.1619) but not statistically significant (p = 0.37), indicating no strong evidence for meaningful change in the relationship between xG differential and points earned across the two seasons. This suggests that the importance of creating higher-quality chances than opponents remained relatively consistent in the transition from Klopp to Slot, while other performance aspects showed more substantial evolution.

My analysis confirmed the significance of home advantage, with matches at Anfield yielding approximately 0.48 additional points on average compared to away fixtures (p $<$ 0.05). This home field advantage remained consistent across both seasons, with no significant interaction between venue and season included in the model. This finding aligns with established research on home advantage in football \citep{kern2012} and highlights the continued importance of venue effects even as other performance dynamics shifted between managerial approaches.

\subsection*{Interpretation}

Interestingly, the base level of performance did not significantly differ between seasons, with the Season main effect showing a non-significant coefficient of -0.1274 (p = 0.565). This suggests that the overall expected points for a match with average xG difference and shot conversion remained relatively stable across the two seasons, despite Liverpool's improved league position from third place to champions.

The most notable seasonal difference emerged in the relationship between shot conversion rate and match outcomes, as indicated by the significant negative interaction term. This finding suggests a more complex dynamic than simply improved finishing efficiency under Slot's management. It points to a potential tactical shift where consistent shot creation quality became more reliably translated into results, reducing the importance of exceptional finishing efficiency for securing points.

Despite this negative interaction term for shot conversion rate, Liverpool's overall performance improved in the 2024-25 season, particularly in terms of home form consistency. This suggests that Slot's approach prioritized consistent performance across different match contexts rather than relying on exceptional finishing in particular matches, resulting in more reliable point accumulation throughout the season.

\subsection*{Limitations to Consider}

While my analysis provides valuable insights, I should acknowledge several limitations. The model's adjusted R-squared of approximately 37\% indicates that substantial outcome variance remains unexplained by the included predictors. Additional factors not captured in my model, such as possession, opposition quality, player availability, or specific tactical adjustments, among others, likely influenced results beyond the metrics examined.

A key methodological challenge is potential endogeneity; that match circumstances may influence team behavior in ways that create reverse causality. For instance, teams that secure an early lead often adopt more defensive approaches, potentially affecting metrics like shot conversion and xG difference within the same match being analyzed. This dynamic relationship between match state and performance metrics represents a limitation in inferring purely causal relationships.

In selecting variables, I deliberately prioritized metrics with more direct relationships to match outcomes. While expected assists, detailed defensive statistics, and possession data could offer additional tactical insight, their relationship to points earned is often mediated through the metrics I've included. High possession or expected assists don't necessarily translate to more points without effective shot conversion, and defensive quality is partially captured in xG difference. Due to the relatively small sample size, I focused on crafting a lean, focused analysis with tangible, heavily weighted metrics, rather than crowding the model with variables that might be difficult to quantify or have less direct impact on results. This offers an exciting area to explore in future research, though.

The limited sample size of matches constrains statistical power for detecting subtle effects. This constraint may have contributed to the non-significance of some parameters, particularly interaction terms that require substantial data to estimate precisely. Further research with larger samples could potentially reveal additional patterns not detectable in my current analysis.

Despite these limitations, the identified relationships between performance metrics and match outcomes provide meaningful insight into the factors that contributed to Liverpool's improved performance under Arne Slot compared to the previous season under Jürgen Klopp.

\section*{Conclusion}

This analysis looked to uncovered the key performance metrics that drove Liverpool's improved results in the 2024-25 Premier League season under Arne Slot. Shot conversion rate showed the strongest effect on points earned, though the significant negative interaction with season revealed this relationship weakened during the championship campaign. This pattern suggests a shift toward greater consistency rather than reliance on exceptional finishing in specific matches, aligning with research on the relationship between expected goals, finishing efficiency, and match outcomes \citep{anzer2023}.

Expected goals difference also significantly predicted points earned, with a smaller but meaningful impact. The non-significant interaction between xG difference and season indicates that chance creation quality remained similarly important under both managers' approaches.

Home advantage consistently delivered approximately half a point more per match at Anfield than away fixtures. While this effect persisted across seasons, Liverpool's championship campaign featured more consistent performance regardless of venue, suggesting improved adaptability under Slot. This finding reinforces previous research on contextual factors in Premier League performance \citep{kern2012}.

Looking ahead, several promising research directions could extend this work, including incorporating a range of additional variables to improve model explanatory power and provide deeper insight into tactical differences between managers. Comparative studies with other successful managerial transitions might also establish whether these patterns represent common features of effective leadership changes or are unique to Liverpool's situation.

My regression model explained approximately 37\% of match outcome variance, highlighting that additional factors beyond those studied likely influenced Liverpool's results. This limitation reflects football performance's inherent complexity and points to opportunities for future research.

For clubs navigating managerial transitions, these findings suggest maintaining consistent performance across different match contexts may prove more valuable than exceptional efficiency in particular fixtures. They highlight the importance of preserving core performance strengths while enhancing aspects like consistent conversion across different match scenarios, potentially providing a template for successful leadership transitions.

Methodologically, this analysis demonstrates the value of regression with interaction terms for examining performance changes across different periods. The variable centering technique I employed offers a useful approach for addressing multicollinearity concerns in sports performance analyses involving interaction effects.

The Liverpool case provides compelling evidence that championship-level performance can be maintained or enhanced during managerial transitions when core strengths are preserved while tactical adjustments improve consistency and adaptability. This research contributes to our understanding of how advanced metrics can highlight subtle shifts in tactical approaches, revealing factors influencing sustained success in elite football that traditional statistics might miss.

\clearpage

\begin{thebibliography}{99}

\bibitem[Anzer and Bauer(2023)]{anzer2023}
Anzer, G., \& Bauer, P. (2023). Expected goals in football: Improving model performance and demonstrating value. \textit{PLOS One}, 8(2), e0282295.

\bibitem[Cavus and Biecek(2022)]{cavus2022}
Cavus, M., \& Biecek, P. (2022). Explainable expected goal models for performance analysis in football analytics. \textit{arXiv:2206.07212}.

\bibitem[Deb and Dey(2019)]{deb2019}
Deb, S., \& Dey, D. (2019). Spatial modeling of shot conversion in soccer to single out goalscoring ability. \textit{Journal of Sports Analytics}, 5(1), 55-64.

\bibitem[Işık(2020)]{isik2020}
Işık, A. A. (2020). Liverpool FC's Performance in Top 6 and Derby Matches - A Statistical Analysis. \textit{ResearchGate}, 342053178.

\bibitem[Jaseziv(2022)]{jaseziv2022}
Jaseziv, J. (2022). worldfootballR: Extract and clean world football (soccer) data. R package version 0.5.8.

\bibitem[Kern et al.(2012)]{kern2012}
Kern, A., Schwarzmann, M., \& Wiedenegger, A. (2012). Measuring the efficiency of English Premier League football: A two-stage data envelopment analysis approach. \textit{Sport, Business and Management}, 2(3), 177-195.

\bibitem[O'Brien(2007)]{obrien2007}
O'Brien, R. M. (2007). A caution regarding rules of thumb for variance inflation factors. \textit{Quality \& Quantity}, 41(5), 673-690.

\bibitem[Roccetti et al.(2024)]{roccetti2024}
Roccetti, M., Berveglieri, F., \& Cappiello, G. (2024). Football Data Analysis: The Predictive Power of Expected Goals (xG). \textit{Proceedings of the 25th International Conference on Intelligent Games and Simulation}, 20-24.

\end{thebibliography}

\clearpage
\section*{Tables}

\begin{table}[h]
\centering
\caption{Regression Analysis of Factors Influencing Liverpool's Points per Match}
\begin{tabular}{lcccc}
\toprule
\textbf{Term} & \textbf{Estimate} & \textbf{Std. Error} & \textbf{t value} & \textbf{Pr($>$\textbar t\textbar)} \\
\midrule
(Intercept) & 2.1853 & 0.1922 & 11.369 & $<$ 2e-16 *** \\
xG\_difference\_c & 0.2148 & 0.1024 & 2.096 & 0.04021 * \\
ShotConversionRate\_c & 11.1988 & 2.1873 & 5.120 & 3.3e-06 *** \\
HomeGame & 0.4803 & 0.2272 & 2.114 & 0.03858 * \\
Season & -0.1274 & 0.2201 & -0.579 & 0.56504 \\
xG\_difference\_c:Season & 0.1619 & 0.1793 & 0.903 & 0.36999 \\
ShotConversionRate\_c:Season & -8.0318 & 2.9506 & -2.722 & 0.00844 ** \\
\bottomrule
\end{tabular}
\begin{tabular}{l}

Significance levels: *** p$<$0.001, ** p$<$0.01, * p$<$0.05, . p$<$0.1 \\
Residual standard error: 0.8667 on 61 degrees of freedom \\
Multiple R-squared: 0.4281, Adjusted R-squared: 0.3718 \\
F-statistic: 7.609 on 6 and 61 DF, p-value: 4.057e-06 \\
\end{tabular}
\end{table}

\begin{table}[h]
\centering
\caption{Variance Inflation Factors}
\begin{tabular}{lc}
\toprule
\textbf{Variable} & \textbf{VIF} \\
\midrule
xG\_difference\_c & 1.6791 \\
ShotConversionRate\_c & 2.4344 \\
HomeGame & 1.1679 \\
Season & 1.0967 \\
xG\_difference\_c:Season & 1.6142 \\
ShotConversionRate\_c:Season & 2.3571 \\
\bottomrule
\end{tabular}
\end{table}

\clearpage
\section*{Figures}

\begin{figure}[H]
\centering
\includegraphics[width=0.8\textwidth]{Figure1.pdf}
\caption{Practical Impact of Key Factors on Points}
\end{figure}

\begin{figure}[H]
\centering
\includegraphics[width=0.8\textwidth]{Figure2.pdf}
\caption{Average Points by Season and Location}
\end{figure}

\begin{figure}[H]
\centering
\includegraphics[width=0.8\textwidth]{Figure3.pdf}
\caption{Coefficient Estimates with 95\% Confidence Intervals}
\end{figure}

\begin{figure}[H]
\centering
\includegraphics[width=0.8\textwidth]{Figure4.pdf}
\caption{Interaction Effect: Shot Conversion Rate by Season}
\end{figure}

\begin{figure}[H]
\centering
\includegraphics[width=0.8\textwidth]{Figure5.pdf}
\caption{Home vs Away Game Effect}
\end{figure}

\begin{figure}[H]
\centering
\includegraphics[width=0.8\textwidth]{Figure6.pdf}
\caption{Observed vs Predicted Points}
\end{figure}

\begin{figure}[H]
\centering
\includegraphics[width=0.8\textwidth]{Figure7.pdf}
\caption{Residual Analysis}
\end{figure}

\end{document}